\documentclass[12pt]{article}

\usepackage{fullpage}
\usepackage{mdframed}
\usepackage{multicol}
\usepackage{colonequals}
\usepackage{algpseudocode}
\usepackage{algorithm}
\usepackage[most, breakable]{tcolorbox}
\usepackage[all]{xy}
\usepackage{proof}
\usepackage{mathtools}
\usepackage{bbm}
\usepackage{amssymb}
\usepackage{amsthm}
\usepackage{amsmath}
\usepackage{amsxtra}
\usepackage{tikz-cd}
\usepackage{enumitem}
\usepackage{euler}
\usepackage{mathpazo}
\usepackage{wrapfig}

\newtheorem{theorem}{Theorem}[section]
\newtheorem{theorem*}{Theorem}
\newtheorem{definition}[theorem]{Definition}
\newtheorem{corollary}{Corollary}[theorem]
\newtheorem{lemma}[theorem]{Lemma}
\newtheorem{lemmact}[theorem]{Lemmact}
\newtheorem{proposition}[theorem]{Proposition}
\newtheorem{remark}[theorem]{Remark}
\newtheorem{example}{Example}[section]


% \newtheorem*{exercisehelper}{Exercise.}
% \newenvironment{exercise}[1]{%
%   \IfBlankTF{#1}
%     {\renewcommand{\exercisehelper}{\textbf{Exercise} \unskip}}
%     {\renewcommand\exercisehelper{\textbf{Exercise #1}}}%
%   \exercisehelper
% }{\endexercisehelper}

\theoremstyle{remark}
\newtheorem{exercise}{Exercise}[section]
\newtheorem*{solution}{Solution}
\newtheorem*{notation}{Notation}
\newcommand{\mathcat}[1]{\textup{\textbf{\textsf{#1}}}} % for defined terms

\newenvironment{problem}[1]
{ \begin{tcolorbox}[breakable]\noindent\textbf{Problem #1}.}
{\vskip 6pt \end{tcolorbox}}

\newenvironment{enumalph}
{\begin{enumerate}\renewcommand{\labelenumi}{\textnormal{(\alph{enumi})}}}
{\end{enumerate}}

\newenvironment{enumroman}
{\begin{enumerate}\renewcommand{\labelenumi}{\textnormal{(\roman{enumi})}}}
{\end{enumerate}}

\newcommand{\defi}[1]{\textsf{#1}} % for defined terms



\setlength{\hfuzz}{4pt}

\let\H\relax
\let\P\relax
\newcommand{\bb}{\mathbb}
\newcommand{\A}{\mathbb A}
\newcommand{\H}{\mathbb H}
\newcommand{\E}{\mathbb E}
\newcommand{\P}{\mathbb P}
\newcommand{\C}{\mathbb C}
\newcommand{\I}{\mathbb I}
\newcommand{\K}{\mathbb K}
\newcommand{\N}{\mathbb N}
\newcommand{\Q}{\mathbb Q}
\newcommand{\R}{\mathbb R}
\newcommand{\Z}{\mathbb Z}
\newcommand{\F}{\mathbb F}
\newcommand{\V}{\mathbb V}
\newcommand{\br}{\mathbf{r}}
\newcommand{\RP}{\mathbb{RP}}
\newcommand{\CP}{\mathbb{CP}}
\newcommand{\nbit}[1]{\{0, 1\}^{#1}}
\newcommand{\bits}{\{0, 1\}^{n}}
\newcommand{\bbni}{\bigbreak \noindent}
\newcommand{\norm}[1]{\left\vert\left\vert#1\right\vert\right\vert}
\newcommand{\dbar}{\overline{\partial}}
\let\d\relax
\newcommand{\d}{\partial}
\newcommand{\e}{\epsilon}
\newcommand{\cO}{\mathcal{O}}
\newcommand{\FF}{\mathcal{F}}
\newcommand{\cG}{\mathcal{G}}
\newcommand{\HH}{\mathcal{H}}
\newcommand{\EE}{\mathcal{E}}
\newcommand{\CC}{\mathcal{C}}
\newcommand{\DD}{\mathcal{D}}
\newcommand{\PP}{\mathcal{P}}
\newcommand{\iso}{\cong}
\newcommand{\tism}{Talked to Sair Shaikh'26}
\DeclareMathOperator{\fm}{\mathfrak{m}}
\DeclareMathOperator{\fn}{\mathfrak{n}}
\DeclareMathOperator{\fp}{\mathfrak{p}}
\DeclareMathOperator{\fq}{\mathfrak{q}}
\DeclareMathOperator{\fM}{\mathfrak{M}}
\DeclareMathOperator{\fP}{\mathfrak{P}}


\newcommand{\Frac}{\text{Frac}}
\newcommand{\Tor}{\text{Tor}}
\newcommand{\Ext}{\text{Ext}}
\newcommand{\gr}{\text{gr}}
\newcommand{\pd}{\text{pd}}
\newcommand{\Proj}{\text{Proj}}
\newcommand{\NF}{\text{NF}}
\newcommand{\Jac}{\text{Jac}}
\newcommand{\red}{\text{red}}
\newcommand{\Nil}{\text{Nil}}
\newcommand{\LC}{\text{LC}}
\newcommand{\LM}{\text{LM}}
\newcommand{\LT}{\text{LT}}
\newcommand{\GL}{\text{GL}}
\newcommand{\tp}{\text{T}}
\newcommand{\init}{\text{in}}
\DeclareMathOperator{\Bl}{Bl}


\DeclareMathOperator{\cS}{\mathscr S}
\DeclareMathOperator{\cV}{\mathscr V}
\DeclareMathOperator{\cR}{\mathscr R}
\DeclareMathOperator{\cT}{\mathscr T}
\DeclareMathOperator{\comRing}{\textbf{comRing}}
\DeclareMathOperator{\Ann}{\text{Ann}}
\DeclareMathOperator{\adj}{\text{Adj}}
\DeclareMathOperator{\Top}{\textbf{Top}}
\DeclareMathOperator{\Mod}{\textbf{Mod}}
% \renewcommand{\Top}{\textbf{Top}}
\DeclareMathOperator{\Set}{\textbf{Set}}
\DeclareMathOperator{\alg}{\textbf{alg}}
\DeclareMathOperator{\ord}{\textbf{ord}}
\DeclareMathOperator{\height}{\text{ht}}
\DeclareMathOperator{\trdeg}{\text{tr.deg}}
\DeclareMathOperator{\std}{\text{std}}
\DeclareMathOperator{\depth}{\text{depth}}
\DeclareMathOperator{\ini}{\text{in}}




\let\1\relax
\newcommand{\1}{\mathbf{1}}
\newcommand{\fr}[2]{\left(\frac{#1}{#2}\right)}
\newcommand{\todo}[1]{\textcolor{red}{\textbf{TODO:} #1}}
\newcommand{\vecz}{\mathbf{z}}
\newcommand{\vecr}{\mathbf{r}}
\DeclareMathOperator{\ot}{\leftarrow}
\DeclareMathOperator{\Cinf}{C^{\infty}}
\DeclareMathOperator{\Id}{Id}
\DeclareMathOperator{\Ell}{Ell}
\DeclareMathOperator{\CL}{\mathcal{CL}}
\DeclareMathOperator{\Alt}{Alt}
\DeclareMathOperator{\Aut}{Aut}
\DeclareMathOperator{\ann}{ann}
\DeclareMathOperator{\codim}{codim}
\DeclareMathOperator{\End}{End}
\DeclareMathOperator{\Hom}{Hom}
\DeclareMathOperator{\id}{id}
\DeclareMathOperator{\im}{im}
\DeclareMathOperator{\M}{M}
\DeclareMathOperator{\Mat}{Mat}
\DeclareMathOperator{\Ob}{Ob}
\DeclareMathOperator{\opchar}{char}
\DeclareMathOperator{\opspan}{span}
\DeclareMathOperator{\rk}{rk}
\DeclareMathOperator{\ev}{ev}
\DeclareMathOperator{\sgn}{sgn}
\DeclareMathOperator{\Sym}{Sym}
\DeclareMathOperator{\tr}{tr}
\DeclareMathOperator{\img}{img}
\DeclareMathOperator{\coker}{coker}
\DeclareMathOperator{\Spec}{Spec}
\DeclareMathOperator{\mSpec}{mSpec}
\DeclareMathOperator{\CandE}{CandE}
\DeclareMathOperator{\CandO}{CandO}
\DeclareMathOperator{\argmax}{argmax}
\DeclareMathOperator{\first}{first}
\DeclareMathOperator{\last}{last}
\DeclareMathOperator{\cost}{cost}
\DeclareMathOperator{\dist}{dist}
\DeclareMathOperator{\path}{path}
\DeclareMathOperator{\parent}{parent}
\DeclareMathOperator{\argmin}{argmin}
\DeclareMathOperator{\excess}{excess}
\let\Pr\relax
\DeclareMathOperator{\Pr}{\mathbf{Pr}}
\DeclareMathOperator{\Exp}{\mathbb{E}}
\DeclareMathOperator{\Var}{\mathbf{Var}}
\let\limsup\relax
\DeclareMathOperator{\limsup}{limsup}
%Paired Delims
\DeclarePairedDelimiter\ceil{\lceil}{\rceil}
\let\oldceil\ceil
\renewcommand{\ceil}[1]{\oldceil*{#1}}

\DeclarePairedDelimiter{\floor}{\lfloor}{\rfloor}
\let\oldfloor\floor
\renewcommand{\floor}[1]{\oldfloor*{#1}}


\newcommand{\circled}[2][]{
  \tikz[baseline=(char.base)]{
    \node[shape=circle, draw, inner sep=1.5pt, #1] (char) {$#2$};
  }
}
\newcommand{\dcircled}[2][]{
  \tikz[baseline=(X.base)]{
    \node[circle, draw, double, double distance=2pt, inner sep=1.5pt, #1] (X) {$#2$};
  }
}

\newcommand{\dagstar}{*}

\newcommand{\tbigwedge}{{\textstyle{\bigwedge}}}
\setlength{\parindent}{0pt}
\setlength{\parskip}{5pt}


\usepackage{listings}
\usepackage{courier}
\usepackage{microtype}


\lstset{
  basicstyle=\footnotesize\ttfamily,
  breaklines=true,
  breakatwhitespace=true
  columns=fullflexible,
  keepspaces=true,
  frame=single,
  escapeinside={(*@}{@*)}
}

\begin{document}

Throughout, let $S = k[x_1, \ldots, x_r]$ with $k$ algebraically closed,  $\fm = (x_1, \ldots, x_r)$, let $J \subseteq \fm$ be an ideal, $R = S/J$, $X = \V(J) \subseteq \A^r$, and $I = \fm R$. The goal of this worksheet is to prove 
\[\gr_I R \;\cong\; S/\init_I(J)\]
and to verify the easy half of the geometric statement on a concrete example.

\begin{problem}{1} \textbf{The Graded Pieces.}

    \begin{enumalph}
        \item Show that $I^n = (\fm^n + J)/J$ as ideals of $R = S/J$. 
        \emph{(Hint: $I = \fm R$ is generated by the images of $x_1, \ldots, x_r$; 
        describe $I^n$ as the image of $\fm^n$ under $S \twoheadrightarrow R$.)}

        \item Using part (a), establish the isomorphism of $k$-vector spaces
        \[\frac{I^n R}{I^{n+1} R} \;\cong\; \frac{\fm^n + J}{\fm^{n+1} + J}.\]
        \emph{(Hint: apply the third isomorphism theorem to the quotient of 
        quotients.)}

        \item The polynomial ring $S$ is graded by degree, $S = \bigoplus_n S_n$ 
        where $S_n$ is the span of degree-$n$ monomials. Show that 
        $S_n \cong \fm^n/\fm^{n+1}$ as $k$-vector spaces, and conclude that 
        $\gr_{\fm} S \cong S$ as graded rings. 
    \end{enumalph}
\end{problem}

% \begin{solution}
%     \bbni
%     \begin{enumalph}
%         \item The ideal $I = \fm R$ is the image of $\fm$ under the quotient $q : S \twoheadrightarrow S/J = R$. Since $q$ is a surjective ring homomorphism, $q(\fm^n) = q(\fm)^n = I^n$. As $q$ has kernel $J$, the preimage of $I^n$ is $\fm^n + J$, so
%         \[I^n = q(\fm^n) = (\fm^n + J)/J.\]

%         \item By (a), $I^n R = (\fm^n + J)/J$ and $I^{n+1}R = (\fm^{n+1}+J)/J$. Since $\fm^{n+1} + J \subseteq \fm^n + J$, the third isomorphism theorem gives
%         \[\frac{I^n R}{I^{n+1}R} = \frac{(\fm^n+J)/J}{(\fm^{n+1}+J)/J} \cong \frac{\fm^n + J}{\fm^{n+1}+J}.\]

%         \item Let $S_n$ be the span of degree-$n$ monomials. Every $g \in \fm^n$ is a polynomial with no terms of degree $< n$; write $g = g_n + g_{n+1} + \cdots$ with $g_i \in S_i$. Modulo $\fm^{n+1}$ (which contains all terms of degree $\geq n+1$) only $g_n$ survives, so the map $g + \fm^{n+1} \mapsto g_n$ is a well-defined $k$-linear isomorphism $\fm^n/\fm^{n+1} \xrightarrow{\sim} S_n$. Summing over $n$ gives a graded $k$-vector space isomorphism $\gr_{\fm} S \cong \bigoplus_n S_n = S$. Finally this respects multiplication: for $g \in S_n, h \in S_m$, the product in $\gr_{\fm} S$ is the class of $gh$ in $\fm^{n+m}/\fm^{n+m+1}$. Since $g$ and $h$ are homogeneous, $gh$ is homogeneous of degree $n+m$, so it maps to $gh \in S_{n+m}$ under the isomorphism above. So $\gr_{\fm} S \cong S$ as graded rings.
%     \end{enumalph}
% \end{solution}

% \newpage 

\begin{problem}{2} \textbf{Initial Forms}
    \begin{enumalph}     
        \item Recall from Section 2 that for $f \in S$ with $f \in \fm^m \setminus  \fm^{m+1}$, the initial form $\init(f) \in \fm^m/\fm^{m+1} = S_m$ is the  lowest-degree homogeneous part of $f$. Define a graded ring  homomorphism 
        \[\phi : S \cong \gr_{\fm} S \longrightarrow \gr_I R, \qquad  S_n \ni \overline{g} \longmapsto \big(g + \fm^{n+1} + J\big)  \in \frac{I^n R}{I^{n+1}R}.\] 
        
        Using Problem 1, show that $\phi$ is surjective.  

        \item Show that for $g \in S_n$ (homogeneous of degree $n$), 
        \[g \in \ker \phi \iff g \in \fm^{n+1} + J\] 
        Then show that $(\ker(\phi))_n = \init_I(J) \cap S_n$ and thus deduce that $\ker\phi = \init_I(J)$, the ideal generated by 
        $\{\init(f) : f \in J\}$.

        \item Conclude from parts (a) and (b) that
        \[\gr_I R \cong S/\init_I(J).\]

        \item \emph{(Why all of $J$ is needed.)} Revisit the Remark in Section 2 with 
        $J = (x^2 + y^3, xy) \subset k[x,y]$. Show that $y^4 \in \init_I(J)$ but 
        $y^4 \notin (\init(x^2+y^3), \init(xy)) = (x^2, xy)$. Explain why this 
        means $\gr_I R$ \emph{cannot} be computed by only taking initial forms 
        of a chosen generating set of $J$. 
    \end{enumalph}
\end{problem}

% \begin{solution}
%     \bbni
%     \begin{enumalph}
%         \item A homogeneous class in $I^nR/I^{n+1}R \cong (\fm^n+J)/(\fm^{n+1}+J)$ is represented by some $a \in \fm^n + J$, say $a = g + j$ with $g \in \fm^n$, $j \in J$. Then $a \equiv g \pmod{\fm^{n+1}+J}$. Writing $g = g_n + g_{n+1} + \cdots$ with $g_i \in S_i$ we have $g \equiv g_n \pmod{\mathfrak{m}^{n+1}}$, so $a \equiv g_n \pmod{\mathfrak{m}^{n+1}+J}$ and the class is $\phi(g_n)$ with $g_n \in S_n$. Hence every homogeneous piece of $\gr_I R$ is hit, and since $\phi$ is graded, $\phi$ is surjective.
        
%         \item We compute $\ker\phi$ degree by degree. For homogeneous $g \in S_n$,
%         \[g \in \ker\phi \iff g \in \fm^{n+1} + J \quad\text{(in $S$).} \tag{$\ast$}\]
%         We claim the homogeneous part of $\ker\phi$ in degree $n$ is exactly $\init_I(J) \cap S_n$.
        
%         {($\supseteq$)} 
        
%         Let $f \in J$ with $\init(f) = f_m \in S_m$, where $f = f_m + f_{m+1} + \cdots$. Then $f_m = f - (f_{m+1} + \cdots) \in J + \fm^{m+1}$, so by $(\ast)$, $f_m \in \ker\phi$. Thus every generator $\init(f)$ of $\init_I(J)$ lies in $\ker\phi$, and since $\ker\phi$ is an ideal, $\init_I(J) \subseteq \ker\phi$.
        
%         {($\subseteq$)} 
        
%         Conversely suppose $g \in S_n$ satisfies $(\ast)$, i.e.\ $g = h + j$ with $h \in \fm^{n+1}$ and $j \in J$. Then $j = g - h$, and since $g \in S_n$ and $h \in \fm^{n+1}$ (all terms of degree $\geq n+1$), the lowest-degree part of $j$ is exactly $g$ provided $g \neq 0$; that is, $\init(j) = g$. Hence $g = \init(j) \in \init_I(J)$. (If $g = 0$ there is nothing to prove.) Combining both inclusions, $(\ker\phi)_n = \init_I(J) \cap S_n$ for every $n$, and summing over $n$ gives $\ker\phi = \init_I(J)$.

%         \item $\phi : S \to \gr_I R$ is a surjective (a) graded ring homomorphism with  $\ker\phi = \init_I(J)$ (b). By the first isomorphism theorem,
%         \[\gr_I R \cong S/\init_I(J) = k[x_1,\ldots,x_r]/\init_I(J).\]

%         \item With $J = (x^2+y^3,\, xy) \subset k[x,y]$ and $\fm = (x,y)$: as shown in the  Remark in Section 2, \[y\cdot(x^2+y^3) - x\cdot(xy) = y^4 \in J,\] and $\init(y^4) = y^4$, so $y^4 \in \init_I(J)$. But every element of  $(\init(x^2+y^3), \init(xy)) = (x^2, xy)$ has $x$ as a factor, while $y^4$  does not, so $y^4 \notin (x^2, xy)$. Thus \[\init_I(J) \supsetneq (x^2, xy) = \big(\init(x^2+y^3), \init(xy)\big).\] Thus, computing $\gr_I R$ by quotienting only by the initial forms  of the chosen generators $x^2+y^3, xy$ would give the \emph{wrong} ring  $k[x,y]/(x^2,xy)$; the correct ring requires the initial forms of all  elements of $J$, including the ``hidden'' generator $y^4$.    
%     \end{enumalph}
% \end{solution}

% \newpage 

\begin{problem}{3}
    \begin{enumalph}
        \item Let $f \in J$ with initial form $\init(f) = f_m$ (its degree-$m$ part, $f_m \neq 0$), so $f = f_m + f_{m+1} + \cdots + f_d$ with each $f_i$ homogeneous of degree $i$. Fix a nonzero direction $p = (p_1, \ldots, p_r)$ and consider the points $tp \in \A^r$ for $t \in k$. Compute $f(tp)$ as a polynomial in $t$ and show \[f(tp) = t^m\big(f_m(p) + t\, f_{m+1}(p) + \cdots + t^{d-m} f_d(p)\big).\]

        \item Suppose $p_n \in X \setminus \{0\}$ is a sequence with $p_n \to 0$ 
        and $[p_n] \to [p]$ in $\P^{r-1}$. Using that $f(p_n) = 0$ for all $n$ 
        (since $p_n \in X$ and $f \in J$), argue informally that $f_m(p) = 0$, 
        i.e.\ the limiting direction $[p]$ lies in $\V(\init(f))$. Conclude 
        \[TC_0(X) \subseteq \V(\init_I(J)).\]
        \emph{(This is the easy inclusion. The reverse inclusion holds when $k$ 
        is algebraically closed but requires tools beyond these notes.)}

        \item \emph{(Worked example.)} For $f = y^2 - x^3$, the cuspidal cubic curve, carry out part (a) explicitly: find $\init(f)$, factor it, and identify the limiting secant direction in $\P^1$.
    \end{enumalph}
\end{problem}

% \begin{solution}
%     \bbni
%     \begin{enumalph}
%         \item Write $f = f_m + f_{m+1} + \cdots + f_d$ with $f_i \in S_i$. Each $f_i$ is  homogeneous of degree $i$, so $f_i(tp) = t^i f_i(p)$. Thus,
%         \[f(tp) = \sum_{i=m}^d t^i f_i(p)  = t^m\big(f_m(p) + t f_{m+1}(p) + \cdots + t^{d-m}f_d(p)\big).\]  
        
%         \item Write each $p_n = t_n v_n$ where $t_n \to 0$ in $k$ and $[v_n] = [p_n] \to [p]$, with representatives $v_n \to v$, $[v] = [p]$. By part (a) applied to  $p_n$, \[0 = f(p_n) = f(t_n v_n)  = t_n^m\big(f_m(v_n) + t_n f_{m+1}(v_n) + \cdots\big).\] Since $t_n \neq 0$ (as $p_n \neq 0$), divide by $t_n^m$: \[f_m(v_n) + t_n f_{m+1}(v_n) + \cdots + t_n^{d-m} f_d(v_n) = 0.\] Let $n \to \infty$: each $f_i$ is a polynomial, hence continuous, so  $f_i(v_n) \to f_i(v)$, and the terms with a factor $t_n^{\geq 1}$ vanish in  the limit. Thus $f_m(v) = 0$, i.e.\ $\init(f)(p) = 0$. Since this holds for  every $f \in J$, the limiting direction $[p]$ lies in $\V(\init_I(J))$,  giving $TC_0(X) \subseteq \V(\init_I(J))$.  
        
%         \item $f = y^2 - x^3$ has a cusp at the origin. The two homogeneous parts are  $y^2$ (degree $2$) and $x^3$ (degree $3$), so the initial form is the  lower-degree one, 
%         \[\init(f) = y^2.\] 
%         Hence the tangent cone is $\V(y^2) \subseteq \A^2$. As a set this is the  single line $\{y = 0\}$ (the $x$-axis), so the cusp has only \emph{one}  limiting secant direction, $[1:0] \in \P^1$ --- in contrast to the node,  which had two. The coordinate ring of the tangent cone is 
%         \[k[x,y]/(y^2) \cong \gr_I R,\] 
%         which is \emph{not reduced}: the factor $y^2$ records that the single  tangent line occurs ``with multiplicity two,'' reflecting how the two  branches of the cusp come together tangentially rather than crossing  transversally. 
%       \end{enumalph}
% \end{solution}


\end{document}